\documentclass{Manual}

\SetProyecto{Revista Universitaria}


% ------------------------------------ %

\begin{document}

    \frontmatter

\Portada
\Indices


% ------------------------------------ %

    \mainmatter


\chapter{Objetivos de la revista}

La existencia de esta revista buscaría cumplir 4 objetivos principales:

\begin{enumerate}
    \item Familiarizar a los alumnos con el proceso de publicación de artículos.
    \item Dar a los alumnos material extra el cual presentar en sus exposiciones.
    \item Aumentar el entusiasmo de los estudiantes al permitirles ver su trabajo en un formato similar al de los formatos académicos.
    \item Generar una biblioteca de los Proyectos Integradores llevados a cabo por los alumnos, para así permitir su futura consulta por nuevas generaciones.
        \subitem Podría considerarse incluir tesinas realizadas por alumnos de sexto y décimo cuatrimestre en sus estadías, ya sea dentro del mismo archivo de la revista o como un apartado extra en el sitio web.
\end{enumerate}



\chapter{Identidad gráfica}

Hasta ahora, las inspiraciones para la revista han sido \href{https://www.pentagram.com/work/scientific-american?rel=sector&rel-id=3}{un rediseño de la revista Scientific American}, la revista \href{https://www.frontiersin.org/journals/sustainability/articles/10.3389/frsus.2026.1818134/pdf}{Frontiers}, \href{https://www.sciencedirect.com/science/article/pii/S1198743X26000650/pdfft?pid=1-s2.0-S1198743X26000650-main.pdf}{Elsevier} y algunas ideas en \href{https://pin.it/3s68QTUbZ}{Pinterest}.

Si bien aún pueden realizarse mejoras, se esperará a recibir retroalimentación de docentes y alumnos para aplicarlas. Asímismo, en un futuro también se aceptarán recomendaciones de alumnos y docentes para mejorar el diseño de la revista.

\section{Paleta de colores}

La paleta de colores utilizada en la revista está inspirada en los colores institucionales de UPSIN. A continuación se listan los colores utilizados en la revista y las áreas en que se usa cada uno:

    \input{fig:paleta_de_colores}

        \definecolor{Primario}{RGB}{9,60,113}
        \definecolor{Secundario}{RGB}{0,160,224}
        \definecolor{Acento}{RGB}{138,222,255}
        \definecolor{Complemento}{RGB}{245,230,80}

\begin{itemize}
    \item Color principal (A) - RGB(9,60,113)
        \subitem Contador de tablas y figuras (Tabla X, Figura X)
        \subitem Color de fondo para encabezado en tablas a una columna
        \subitem Enlaces externos
    \item Color secundario (B) - RGB(0,160,224)
        \subitem Enlaces internos (citas, figuras y tablas)
    \item Color de acento (C) - RGB(138,222,255)
    \item Color de complemento (D) - RGB(245,230,80)
    \item Gris (E) - RGB(242,242,242)
        \subitem Color de fondo para filas en tablas (intercalado con blanco)
\end{itemize}

\section{Tipografías}

Se estarían utilizando 2 tipografías: Red Hat Display y Crimson Pro. Cada una sería utilizada en diferentes  áreas.

La tipografía Red Hat Display (\autoref{fig:tipografía_red_hat_display}) sería utilizada para titulos y encabezados, variando y combinando su estilo y peso (negrita, normal e itálica) para cada nivel de titulo o encabezado. También se estaría utilizando para los contadores de figuras y tablas (Tabla X, Figura X).

    \input{fig:tipografía_red_hat_display}

Por otro lado, la tipografía Crimson Pro (\autoref{fig:tipografía_crimson_pro}) sería utilizada para el texto principal, así como descripciones de tablas y figuras.

    \input{fig:tipografía_crimson_pro}


Para el texto general se utilizará un tamaño de 12pt. Para los encabezados de distintos niveles se utilizarán los siguientes tamaños:

\begin{itemize}
    \item 18pt en negritas 
    \item 16pt en negritas
    \item 14pt en itálica
\end{itemize}

\section{Presentación de tablas y figuras}

\subsection{Tablas}

Los encabezados de las tablas tendrán Red Had Display como tipografía, mientras que el resto de la tabla (cuerpo y notas a pie) mantendrán Crimson Pro como tipografía.

El color de las filas del cuerpo alternará entre blanco y gris. En el caso de los encabezados estos podrán utilizar el color primario (azul oscuro, \autoref{tab:ejemplo_tabla_azul}) o blanco como color de fondo (\autoref{tab:ejemplo_tabla_blanco}).

    \input{tablas_ejemplo}

\subsection{Figuras}

En el caso de las figuras se respetará el diseño de los autores originales. Se usará una linea azul oscuro para separar la figura de la descripción.

\begin{figure}[H]
    \centering
    \includegraphics[width=0.85\linewidth]{ejemplo_figura}
    \textcolor{Primario}{\rule{\linewidth}{0.2pt}}
    \caption{Ejemplo de figura}
    \label{fig:ejemplo_figura}
\end{figure}

\chapter{Publicación}

\section{Periodos de publicación}

Dado que el principal objetivo de la revista es dar una mayor difusión de los trabajos realizados por los alumnos en la materia de Proyecto Integrador, la publicación de números nuevos se daría en base a la periodicidad de esta materia.

Por la organización del calendario escolar, esta materia tiene lugar en los periodos de Enero-Abril con 1 grupo, y en Mayo-Agosto con 2 grupos a la vez. Así pues, la publicación de nuevos números se realizaría en Abril y Agosto.

Dado que dos grupos presentarían proyectos en la publicación de Agosto, estaría a consideración si se publica un solo número con los trabajos de ambos grupos, o se divide en una ublicación por grupo. Asímismo, lo anterior se basa en que la revista se enfoque casi exclusivamente en los productos de la asignatura Proyecto Integrador, sin embargo, podrían existir más publicaciones si presenta el suficiente interés por parte de la comunidad estudiantil.

\section{Formatos de publicación}

La publicación de la revista sería en su mayor parte digital. Para esto se crearía una \textbf{\href{https://darahraccoon.github.io/BT_Magazine/}{página web}}\footnote{El diseño de la página aún está en proceso} en la cual se alojarían los números de la revista, así como los requisitos de envío.

Además de esto, se podría considerar realizar una impresión pequeña (3-4 ejemplares) para utilizarlas como recurso extra en las presentaciones de los proyectos.

\section{Tipos de publicaciones}

La revista tendría 2 tipos principales de artículos: Artículo de investigación y reviews.

Este tipo de trabajos son aquellos que los alumnos pueden llevar a cabo en la materia de Proyecto Integrador, por lo que solo sería necesario que lo especifiquen al momento de enviar su artículo.

Se propone para el primer número añadir publicaciones tipo Especial\footnote{El nombre puede cambiar}. Este se propone como un tipo de artículo que expanda la participación de los alumnos en la revista. Consistiría en artículos donde su formato sería más libre, aceptando artículos de divulgación, ensayos, producciones artísticas (poemas e ilustraciones, por ejemplo), etc., y donde las temáticas sean más variadas, tocando áreas de la biotecnología menos exploradas en la universidad, tales como las ciencias sociales, por ejemplo.

Los artículos de esta área los podrían enviar alumnos de cualquier cuatrimestre de manera voluntaria. Ya que estos trabajos no formarían parte de alguna asignatura, su proceso de revisión se establecería dependiendo del interés de los alumnos en publicar y en la disponibilidad de personas revisoras (docentes y/o alumnos). Asímismo, dependiendo del interés de los alumnos podría quedar como una sola categoría o segmentarse en varias.

\section{Proceso de envío}

Para familiarizar a los alumnos con el proceso de envío de artículos a una revista, se podría establecer un formato de envio el cual los alumnos deberán seguir para la publicación de sus artículos. El formato podría ser determinado en base a recomendaciones de docentes para que sea lo más similar posible al estándar académico.

A continuación se presentan algunos de los aspectos a definir en el formato, así como recomendaciones:

\begin{itemize}
    \item Tipo de archivo
        \subitem Google Docs o Word
    \item Características de la hoja
        \subitem Tamaño carta
        \subitem Márgenes de ~2.5cm en los 4 lados (por defecto en programas como Google Docs o Word)
        \subitem Sin encabezado ni pie de página
    \item Características de la tipografía
        \subitem Fuentes: Times New Roman o Arial
        \subitem Tamaño de fuente: 12pt
        \subitem Interlineado de 1.5
        \subitem Titulos a 16pt y en negritas
        \subitem Nombres científicos en itálica
    \item Figuras
        \subitem Enviar siempre archivos de imagen aparte
    \item Tablas
        \subitem De realizarse en hojas de cálculo (Excel, Sheets, etc.), adjuntar tales archivos
    \item Bibliografía
        \subitem Formato APA7
\end{itemize}


Tal como hacen algunas editoriales, se puede establecer también una plantilla que cumpla con los estándares de envío para facilitar la configuración del documento. A través de \textbf{\href{https://docs.google.com/document/d/1M7Ag4gmHuryCKiGC0duv6faO6SFfbrNe-7jA6rKuEZ4/edit?usp=sharing}{este enlace}} se presenta una propuesta de plantilla. 

\section{Revisión de artículos}

Dado que los artículos que se publicarán son producto de una materia con docente asignado, tal docente llevaría a cabo la revisión de los textos.

Aún así, podría existir un proceso de revisión por pares realizado por alumnos de mayor grado. Tal ejercicio podría expandir los beneficios académicos y pedagógicos de la revista al acercar a más alumnos al mundo de la academia y sus procesos.

\section{Proceso de maquetación}

La plantilla de LaTeX mediante la cual se maqueta la revista permite preparar los número unicamente copiando y pegando lo escrito por los alumnos en lugares determinados, por lo que la preparación de la revista sería algo sencillo que podría realizarse por una persona. Personalmente disfruto del proceso, por lo que me gustaría ser el encargado de tal área, sin embargo, dado que LaTeX es una herramienta útil para el procesamiento de archivos, estaría abierto a ceder tal proceso a otros alumnos si estos se muestran interesados a aprender sobre el programa. 



% ------------------------------------ %
    
    \appendix
    


\end{document}